% !TEX program = pdflatex
\documentclass[12pt,a4paper]{article}
\usepackage{amsmath,amssymb,amsthm}
\usepackage{graphicx}
\usepackage[hidelinks]{hyperref}
\hypersetup{pdftitle={Validating the Arb backend of the MST solver}}
\usepackage{xcolor}
\usepackage{booktabs}
\usepackage[margin=2.5cm]{geometry}

\graphicspath{{../../figures/}}

% NOTE: run  julia --project=. examples/arb_validation/arb_amplitude_validation.jl
% to (re)generate the figures before compiling this document.

\title{Validating the Arb (ball-arithmetic) backend of the MST solver\\
\large $B^{\mathrm{inc}}$, $B^{\mathrm{ref}}$, $R^{\mathrm{in}}$, $R^{\mathrm{up}}$
at large complex frequency}
\author{}
\date{\today}

\newcommand{\Binc}{B^{\mathrm{inc}}}
\newcommand{\Bref}{B^{\mathrm{ref}}}
\newcommand{\Rin}{R^{\mathrm{in}}}
\newcommand{\Rup}{R^{\mathrm{up}}}
\newcommand{\eps}{\epsilon}
\newcommand{\Acb}{\texttt{Acb}}
\newcommand{\Arb}{\texttt{Arb}}

\begin{document}
\maketitle

\begin{abstract}
We validate the arbitrary-precision \emph{ball-arithmetic} (Arb) backend of our
Julia MST (Mano--Suzuki--Takasugi) solver for the homogeneous Teukolsky equation,
computing the incidence/reflection amplitudes $\Binc,\Bref$ and the radial
solutions $\Rin,\Rup$ at large complex frequency
$\omega=|\omega|\,e^{i\theta}$, with $|\omega|\in\{2,10\}$ and
$\theta\in[0^\circ,90^\circ)$.  Because Mathematica's \textsf{Teukolsky} package
returns \texttt{Indeterminate} for these quantities at any complex $\omega$, the
requested 100-digit cross-check is infeasible there; we instead validate by
\emph{self-consistency} against the library's independent \textsf{BigFloat}
(MPFR) backend at matched precision, supplemented by the reference-free Wronskian
identity for the radial pair.  The Arb amplitudes agree with the BigFloat
reference to $50$--$74$ digits across the whole frequency sweep, with isolated
ill-conditioned angles whose accuracy is recovered \emph{linearly} by raising the
working precision --- the signature of a correct but ill-conditioned computation.
\end{abstract}

\section{The MST quantities}

The radial Teukolsky function $R_{s\ell m\omega}(r)$ obeys
\begin{equation}
  \Delta^{-s}\frac{d}{dr}\!\Big(\Delta^{s+1}\frac{dR}{dr}\Big)
  +\Big(\frac{K^2-2is(r-M)K}{\Delta}+4is\omega r-\lambda\Big)R=0,
  \qquad \Delta=r^2-2Mr+a^2,
\end{equation}
with $K=(r^2+a^2)\omega-ma$ and $\lambda$ the spin-weighted spheroidal eigenvalue.
The MST method expands the two physical solutions in series of hypergeometric
($\Rin$, near horizon) and Coulomb-wave ($\Rup$, near infinity) functions whose
coefficients $f^\nu_n$ satisfy a three-term recurrence,
\begin{equation}
  \alpha^\nu_n f^\nu_{n+1}+\beta^\nu_n f^\nu_n+\gamma^\nu_n f^\nu_{n-1}=0,
  \label{eq:rec}
\end{equation}
solved by the minimal solution obtained from the two tail continued fractions
$R_n=f_n/f_{n-1}$, $L_n=f_n/f_{n+1}$.  The series converges only for the
\emph{renormalized angular momentum} $\nu$, fixed here from the monodromy of the
confluent Heun equation,
\begin{equation}
  \cos(2\pi\nu)=\cos\!\big(\pi(\mu_1-\mu_2)\big)
  +\frac{2\pi^2}{\Sigma_1\,\Sigma_2}\,(-1)^{n-1}a^{(1)}_{n+1}a^{(2)}_{n+1},
  \qquad
  \nu=\ell-\frac{\arccos\cos(2\pi\nu)}{2\pi},
  \label{eq:mono}
\end{equation}
the $\Sigma_i$ being truncated sums of the monodromy recurrence
$a^{(i)}_n$.  The asymptotic amplitudes follow from the matching coefficients
$K_\nu$ and the factors $A^\nu_\pm$ (Sasaki--Tagoshi 2003, Eqs.~157--170):
\begin{align}
  \Binc&=\omega^{-1}\Big(K_\nu-i\,e^{-i\pi\nu}
        \tfrac{\sin\pi(\nu-s+i\eps)}{\sin\pi(\nu+s-i\eps)}K_{-\nu-1}\Big)
        A^\nu_+\,e^{-i(\eps\ln\eps-\frac{1-\kappa}{2}\eps)},\\
  \Bref&=\omega^{-1-2s}\big(K_\nu+i\,e^{i\pi\nu}K_{-\nu-1}\big)
        A^\nu_-\,e^{+i(\eps\ln\eps-\frac{1-\kappa}{2}\eps)},
\end{align}
with $\eps=2M\omega$, normalised by $B^{\mathrm{trans}}$.  Any two homogeneous
solutions satisfy the $r$-independent Wronskian
\begin{equation}
  W=\Delta^{s+1}\big(\Rin\,(\Rup)'-(\Rin)'\,\Rup\big)=\text{const},
  \label{eq:wron}
\end{equation}
which gives a \emph{reference-free} test: $|W(r_1)-W(r_2)|/|W|$ measures how well
the truncated series actually solves the equation.

\section{The Arb backend and the choice of reference}

The solver is generic over the floating type: \texttt{Float64},
\texttt{BigFloat} (MPFR), and Arb's ball types \Arb/\Acb\ (via
\texttt{Arblib.jl}).  Ball arithmetic carries a rigorous error radius with every
value.  Two facts shape the validation.

\paragraph{No external reference at complex $\omega$.}
Mathematica's \textsf{Teukolsky} package returns \texttt{Indeterminate} for
$\Binc,\Bref,\Rin,\Rup$ whenever $\operatorname{Im}\omega\neq0$ (it resolves only
$\nu$, to $\sim15$ digits).  The literal ``compare against Mathematica at $100$
digits'' is therefore impossible for the complex frequencies of interest.  We
adopt instead a self-consistency reference: the library's BigFloat backend at the
\emph{same} working precision.  Arb and BigFloat share no arithmetic kernels, so
their agreement to working precision certifies the Arb code path (type dispatch,
the special-function bridges, the continued-fraction and monodromy kernels).

\paragraph{Midpoint philosophy.}
At large $|\omega|$ the MST closed forms~\eqref{eq:mono} suffer heavy
cancellation.  Propagated through ball arithmetic, the \emph{certified radius}
grows until it can exceed the midpoint magnitude, and the rigorous comparisons
that drive convergence tests become undecidable.  Following the existing
monodromy driver, we treat converged values as point estimates: ball radii are
stripped at the decision points, and the recurrences~\eqref{eq:rec},
\eqref{eq:mono} are run in \emph{point arithmetic} (midpoints computed exactly to
the working precision, exactly as MPFR), so the Arb midpoints track BigFloat.
The bridges required for \texttt{Complex\{Arb\}} to traverse the generic stack are:
$\log$, $\arccos$ (with a NaN guard for $|z|\gtrsim10^{40}$), $\sqrt{\cdot}$,
\texttt{complex(::Acb)} for the \texttt{Arb/Complex\{Arb\}}$\to$\Acb\ promotion,
and \texttt{loggamma} for the radial ${}_2F_1$; plus a point-arithmetic
continued-fraction and monodromy recurrence.  All are strict no-ops for
\texttt{Float64}/\texttt{BigFloat}.

\section{Results}

\paragraph{Amplitude agreement across the $\theta$ sweep.}
Figures~\ref{fig:w2} and~\ref{fig:w10} show the number of digits to which the Arb
$\Binc,\Bref$ agree with the BigFloat reference at $256$-bit precision, as
$\theta$ runs from $0^\circ$ to $89.99^\circ$.  At $|\omega|=2$ the agreement is a
flat $50$--$74$ digits for both Schwarzschild and Kerr.  At $|\omega|=10$ most
angles still agree to $60$--$74$ digits; a few ill-conditioned angles drop lower
(e.g.\ Kerr $a=0.9$, $\theta=60^\circ$: $\sim7$ digits).

\begin{figure}[t]
  \centering
  \includegraphics[width=0.78\textwidth]{arb_amp_agreement_w2.pdf}
  \caption{Arb-vs-BigFloat agreement (digits) for $\Binc,\Bref$ at $|\omega|=2$,
  $s=-2$, $\ell=m=2$, $256$-bit.  Flat $50$--$74$ digit agreement over the whole
  $\theta$ sweep.}
  \label{fig:w2}
\end{figure}

\begin{figure}[t]
  \centering
  \includegraphics[width=0.78\textwidth]{arb_amp_agreement_w10.pdf}
  \caption{As Fig.~\ref{fig:w2} but at $|\omega|=10$.  Most angles agree to
  $60$--$74$ digits; isolated angles are ill-conditioned (see
  Fig.~\ref{fig:recovery}).}
  \label{fig:w10}
\end{figure}

\paragraph{Precision recovery.}
The low-digit angles are not errors: the cancellation costs a \emph{constant}
number of bits, independent of the working precision.  Figure~\ref{fig:recovery}
shows that for the worst case (Kerr $\theta=60^\circ$, $|\omega|=10$) the
agreement grows linearly --- $7\to84\to238$ digits at $256\to512\to1024$ bits ---
a fixed $\sim230$-bit conditioning loss.  The Schwarzschild $\theta=45^\circ$ case
loses only $\sim25$ bits.  This linear recovery is the signature of a correct but
ill-conditioned computation: to obtain $D$ correct digits at a given mode, use a
working precision exceeding the conditioning loss by $\sim3.3D$ bits.

\begin{figure}[t]
  \centering
  \includegraphics[width=0.72\textwidth]{arb_amp_precision_recovery.pdf}
  \caption{$\Binc$ agreement vs working precision at $|\omega|=10$.  The constant
  vertical gap to the diagonal (full precision) is the conditioning loss: small
  for Schwarzschild $\theta=45^\circ$, $\sim230$ bits for Kerr $\theta=60^\circ$.}
  \label{fig:recovery}
\end{figure}

\paragraph{Radial Wronskian.}
Figure~\ref{fig:wron} tests the radial pair via~\eqref{eq:wron}.  At $|\omega|=0.5$
the Arb solutions are self-consistent to $\gtrsim50$ digits once the radial
truncation $n_{\max}$ is large enough.  At $|\omega|\geq2$ the radial series needs
$n_{\max}\!\propto\!|\omega|$ to converge; the BigFloat path then reaches
$10^{-22}$ ($n_{\max}=240$) and $10^{-46}$ ($n_{\max}=320$), confirming the
\emph{formalism} converges, whereas the Arb path plateaus near $10^{11}$ because
the external hypergeometric (${}_2F_1$) evaluation is carried out in ball
arithmetic and is conditioning-limited there.  Point-arithmetic ${}_2F_1$ for the
Arb radial functions at $|\omega|\geq2$ is the natural next step.

\begin{figure}[t]
  \centering
  \includegraphics[width=0.80\textwidth]{arb_radial_wronskian_nmax.pdf}
  \caption{Radial Wronskian self-consistency $|W(8)-W(20)|/|W|$ vs $n_{\max}$
  (256-bit).  Solid: Arb; dashed: BigFloat.  Arb converges at $|\omega|=0.5$; at
  $|\omega|=2$ the BigFloat formalism converges while the Arb ${}_2F_1$ path is
  conditioning-limited.}
  \label{fig:wron}
\end{figure}

\section{Conclusion}

The Arb backend computes $\Binc,\Bref$ correctly at large complex frequency,
matching the independent BigFloat path to $50$--$74$ digits across
$|\omega|\in\{2,10\}$ and the full $\theta$ sweep, with the only shortfalls being
isolated ill-conditioned angles whose accuracy is fully recovered by raising the
working precision.  The radial solutions $\Rin,\Rup$ are validated at low
$|\omega|$ via the reference-free Wronskian; at $|\omega|\geq2$ they are limited
by the ball-arithmetic hypergeometric evaluation, not by the MST formalism, which
the BigFloat reference confirms still converges.  All Arb-specific code is purely
additive: the \texttt{Float64}/\texttt{BigFloat} paths are byte-for-byte
unchanged and the full test suite passes.

\end{document}
